\section*{Capítulo \ref{ch:g11:hormones-and-reproduction} --- O controle hormonal da reprodução}

\begin{solution}{exo:g11:hormones-and-reproduction:1}
Um \omterm{def:g11:hormones-and-reproduction:hormone}{hormônio} é uma molécula secretada no sangue por uma glândula e que
age sobre \omterm{def:g10:cells-common-unit:cell}{células} distantes que carregam o receptor dele. A
\omterm{def:g11:hormones-and-reproduction:hormone}{retroalimentação} é a ação do produto final sobre as glândulas do começo
da cadeia: negativa quando reduz o comando (mantendo um nível
constante), positiva quando o aumenta (produzindo um pico).
\end{solution}

\begin{solution}{exo:g11:hormones-and-reproduction:2}
Hipotálamo (GnRH), hipófise (LH e FSH), gônada (\omterm{prop:g11:becoming-male-female:hormones}{testosterona} no
testículo; estrogênios e progesterona no ovário).
\end{solution}

\begin{solution}{exo:g11:hormones-and-reproduction:3}
Os túbulos seminíferos produzem os espermatozoides; as \omterm{def:g10:cells-common-unit:cell}{células} entre
eles secretam a \omterm{prop:g11:becoming-male-female:hormones}{testosterona}.
\end{solution}

\begin{solution}{exo:g11:hormones-and-reproduction:4}
Fase folicular (dias 1--13): folículo em crescimento, estrogênios.
Ovulação (dia 14): pico de LH. Fase lútea (dias 15--28): corpo lúteo,
progesterona e estrogênios.
\end{solution}

\begin{solution}{exo:g11:hormones-and-reproduction:5}
A degeneração do corpo lúteo no fim do ciclo: a queda da progesterona e
dos estrogênios faz o revestimento espessado do útero se desfazer e
sangrar.
\end{solution}

\begin{solution}{exo:g11:hormones-and-reproduction:6}
Pico de estrogênio no dia 12, pico de LH no dia 13, ovulação no dia 14,
pico de progesterona no dia 21.
\end{solution}

\begin{solution}{exo:g11:hormones-and-reproduction:7}
A castração retira a \omterm{prop:g11:becoming-male-female:hormones}{testosterona} e, com ela, a inibição do hipotálamo
e da hipófise: o LH sobe. A \omterm{prop:g11:becoming-male-female:hormones}{testosterona} injetada restaura a inibição e
o LH cai --- \omterm{def:g11:hormones-and-reproduction:hormone}{retroalimentação} negativa nos dois sentidos.
\end{solution}

\begin{solution}{exo:g11:hormones-and-reproduction:8}
No dia 12 o nível de estrogênio está alto e sustentado, o que inverte a
\omterm{def:g11:hormones-and-reproduction:hormone}{retroalimentação} para positiva e desencadeia o pico. No dia 20 os
estrogênios estão moderados e a progesterona alta, e a \omterm{def:g11:hormones-and-reproduction:hormone}{retroalimentação}
é negativa: o LH fica contido.
\end{solution}

\begin{solution}{exo:g11:hormones-and-reproduction:9}
Dia 21: a fase lútea dura cerca de 14 dias seja qual for a duração do
ciclo, de modo que a ovulação fica 14 dias antes da menstruação
seguinte; a semana a mais está na fase folicular.
\end{solution}

\begin{solution}{exo:g11:hormones-and-reproduction:10}
O estrogênio constante e moderado da pílula e o \omterm{def:g11:hormones-and-reproduction:hormone}{hormônio} semelhante à
progesterona mantêm a \omterm{def:g11:hormones-and-reproduction:hormone}{retroalimentação} negativa o mês inteiro: o FSH
fica baixo demais para um folículo amadurecer, de modo que não há pico
de estrogênio, nem \omterm{def:g11:hormones-and-reproduction:hormone}{retroalimentação} positiva, nem pico de LH, nem
ovulação.
\end{solution}

\begin{solution}{exo:g11:hormones-and-reproduction:11}
A \omterm{prop:g11:becoming-male-female:hormones}{testosterona} externa inibe o hipotálamo e a hipófise dele; o LH e o
FSH caem; os testículos dele, sem estímulo, param de secretar e param
de fabricar espermatozoides, e encolhem.
\end{solution}

\begin{solution}{exo:g11:hormones-and-reproduction:12}
Sem a inibição do estrogênio, o GnRH e o FSH sobem; os folículos
amadurecem; o pico de estrogênio deles (agindo sobre a hipófise, que
ainda consegue responder) desencadeia um pico de LH e a ovulação. Ele é
usado quando a ovulação falha porque os folículos recebem FSH de menos.
\end{solution}

\begin{solution}{exo:g11:hormones-and-reproduction:13}
O folículo cresce e secreta (o ovário e o FSH funcionam), mas o pico de
estrogênio não produz um pico: a \omterm{def:g11:hormones-and-reproduction:hormone}{retroalimentação} positiva falha no
hipotálamo ou na hipófise. Uma injeção de um \omterm{def:g11:hormones-and-reproduction:hormone}{hormônio} semelhante ao LH
no momento do pico pode desencadear a ovulação.
\end{solution}

\begin{solution}{exo:g11:hormones-and-reproduction:14}
Mantido vivo, o corpo lúteo continua secretando progesterona: o
revestimento é mantido (nenhuma menstruação) e a \omterm{def:g11:hormones-and-reproduction:hormone}{retroalimentação}
negativa segura o LH e o FSH (nenhum folículo novo nem ovulação). O
\omterm{def:g11:hormones-and-reproduction:hormone}{hormônio} aparece no sangue e na urina poucos dias depois da
implantação, e é isso que os testes detectam.
\end{solution}

\begin{solution}{exo:g11:hormones-and-reproduction:15}
$12 \times 0.6 \times 0.4 \approx 2.9$: cerca de 3 embriões.
Probabilidade de nenhum nascimento em três transferências:
$0.65^3 \approx 0.27$; pelo menos um nascimento: cerca de 73\%.
\end{solution}

\begin{solution}{pb:g11:hormones-and-reproduction:1}
\textbf{1.} Dia 14: o pico de LH do dia 13 antecede a ovulação em um
dia, e a progesterona, que só aparece depois da ovulação, sobe a partir
do dia 14.

\textbf{2.} Fase folicular dias 1--13 (13 dias); fase lútea dias 14--27
ou 28 (cerca de 14 dias).

\textbf{3.} O folículo em crescimento no dia 10; o corpo lúteo no dia
22, tanto para os estrogênios quanto para a progesterona.

\textbf{4.} Dia 6: fino, começando a engrossar sob os estrogênios. Dia
22: espesso e secretor sob a progesterona, pronto para um embrião. Dia
28: se desfazendo e sangrando, já que os dois \omterm{def:g11:hormones-and-reproduction:hormone}{hormônios} caíram.

\textbf{5.} No começo, os \omterm{def:g11:hormones-and-reproduction:hormone}{hormônios} do ciclo anterior caíram e a
inibição deles junto, de modo que o FSH sobe; na fase lútea a
progesterona e os estrogênios o inibem fortemente.

\textbf{6.} \omterm{def:g11:hormones-and-reproduction:hormone}{Retroalimentação} negativa: os estrogênios crescentes do
folículo inibem a hipófise, e o FSH cai.

\textbf{7.} \omterm{def:g11:hormones-and-reproduction:hormone}{Retroalimentação} positiva: os estrogênios no máximo deles
por cerca de dois dias estimulam o hipotálamo e a hipófise em vez de
inibi-los, e o LH dispara.

\textbf{8.} Negativa: a progesterona com estrogênios moderados inibe a
hipófise; sem um pico sustentado de estrogênio, não há
\omterm{def:g11:hormones-and-reproduction:hormone}{retroalimentação} positiva nem segundo pico.

\textbf{9.} A progesterona e os estrogênios desabam; o revestimento se
desfaz e sangra no dia 1 do ciclo seguinte; o FSH, liberado da
inibição, sobe.

\textbf{10.} A queda dos \omterm{def:g11:hormones-and-reproduction:hormone}{hormônios} lúteos retira a \omterm{def:g11:hormones-and-reproduction:hormone}{retroalimentação}
negativa, o FSH sobe e um novo folículo começa a crescer: o fim de um
ciclo é o sinal do seguinte.

\textbf{11.} FSH e LH baixos o mês inteiro (\omterm{def:g11:hormones-and-reproduction:hormone}{retroalimentação} negativa
mantida pelos \omterm{def:g11:hormones-and-reproduction:hormone}{hormônios} da pílula); nenhum folículo amadurece.

\textbf{12.} Nenhum pico (nenhum folículo amadurecendo), nenhum pico de
LH, nenhuma ovulação.

\textbf{13.} Os \omterm{def:g11:hormones-and-reproduction:hormone}{hormônios} da pílula caem e o revestimento, fino, se
desfaz e sangra; é um sangramento de privação, e não o fim de uma fase
lútea --- não houve corpo lúteo.

\textbf{14.} Sem os \omterm{def:g11:hormones-and-reproduction:hormone}{hormônios} da pílula a \omterm{def:g11:hormones-and-reproduction:hormone}{retroalimentação} se afrouxa,
o FSH sobe, um folículo pode começar a crescer e a secretar; se ele
chegar ao pico de estrogênio dele, um pico de LH e a ovulação podem vir
depois, e o revestimento fino e o muco espesso deixam de ser mantidos.

\textbf{15.} A progesterona sozinha inibe o hipotálamo e a hipófise o
bastante para impedir o pico de LH, e engrossa o muco do colo do útero.

\textbf{16.} No hipotálamo ou na hipófise: o ovário não é estimulado.
Pulsos de GnRH (se a hipófise funcionar) ou injeções de FSH e LH
restauram o crescimento dos folículos e a ovulação, contornando o nível
que falta.

\textbf{17.} A fecundação in vitro: óvulos recolhidos do ovário,
fecundados numa placa, embrião colocado no útero --- o encontro dos
gametas e o transporte que a tuba uterina teria feito.

\textbf{18.} O pico natural de LH; a ovulação vem cerca de 36 horas
depois do pico, de modo que os óvulos são recolhidos pouco antes de
serem liberados e perdidos.

\textbf{19.} Naturalmente, os estrogênios crescentes do folículo
dominante reduzem o FSH, e os folículos menores, privados de FSH,
degeneram; só o maior sobrevive. O FSH injetado atropela essa seleção e
mantém todos eles crescendo.

\textbf{20.} Ovulação no dia 14; a \omterm{def:g11:hormones-and-reproduction:hormone}{retroalimentação} negativa que
reduziu o FSH à medida que o folículo crescia, depois a
\omterm{def:g11:hormones-and-reproduction:hormone}{retroalimentação} positiva que transformou o pico de estrogênio num pico
de LH; a pílula manteve a \omterm{def:g11:hormones-and-reproduction:hormone}{retroalimentação} negativa o mês inteiro, de
modo que nenhum folículo amadureceu e nenhum pico aconteceu.
\end{solution}
